\begin{table}[H]
\centering
\caption{Effect of a 20.0 Percent Shock to Crude Oil and Gas Prices, 2023}
\scriptsize
\setlength{\tabcolsep}{4.2pt}
\begin{tabular}{lrrrrr}
\toprule
Country & \shortstack{Average inflation\\Total} & \shortstack{Average inflation\\Direct} & \shortstack{Income inflation\\gap} & \shortstack{Direct income\\inflation gap} & \shortstack{Indirect income\\inflation gap}
 \\
\midrule
DE & 1.2 & 0.3 & 0.1 & 0.1 & 0.0
 \\
FR & 1.3 & 0.0 & -0.0 & 0.0 & -0.0
 \\
IT & 2.1 & 0.7 & -0.2 & -0.1 & -0.1
 \\
ES & 1.7 & 0.0 & 0.2 & 0.0 & 0.2
 \\
NL & 2.1 & 0.5 & -0.1 & 0.1 & -0.2
 \\
BE & 1.6 & 0.0 & -0.2 & 0.0 & -0.2
 \\
EA20 & 1.5 & 0.3 & 0.0 & 0.0 & -0.0
 \\
\bottomrule
\end{tabular}
\par\vspace{0.35em}\footnotesize\noindent Sources: Eurostat HICP, HBS, FIGARO, author's calculations.
\par\vspace{0.25em}\footnotesize\noindent Notes. Entries are percentage-point effects of the sectoral input price shock shown in the table title. The price response is computed as $\Delta p = (I - A^{\prime})^{-1}s$, with $A$ built from FIGARO. The direct effect is the component from the shocked sector in the household consumption basket. Indirect is the Leontief propagation component. Average inflation total includes direct and indirect effects before rounding. Income inflation gap = Q1 - Q5 of revenues.
\end{table}
\vspace{0.6em}
